% 2013 辽宁卷文科第 13 题：某几何体的三视图。
\begin{tikzpicture}[font=\footnotesize, line width=0.5pt, >=Stealth]
  % 正视图：宽 4、高 4 的矩形，内部两条虚线分出其宽为 1、2、1
  \begin{scope}[shift={(0,0)}]
    \draw (0,0) rectangle (4,4);
    \draw[dashed] (1,0) -- (1,4);
    \draw[dashed] (3,0) -- (3,4);
    % 上方尺寸标注
    \draw[<->] (0,4.45) -- (1,4.45) node[midway, above] {$1$};
    \draw[<->] (1,4.45) -- (3,4.45) node[midway, above] {$2$};
    \draw[<->] (3,4.45) -- (4,4.45) node[midway, above] {$1$};
    \draw (0,4) -- (0,4.6);
    \draw (1,4) -- (1,4.6);
    \draw (3,4) -- (3,4.6);
    \draw (4,4) -- (4,4.6);
    % 左侧高度标注
    \draw[<->] (-0.45,0) -- (-0.45,4) node[midway, left] {$4$};
    \draw (0,0) -- (-0.6,0);
    \draw (0,4) -- (-0.6,4);
  \end{scope}
  % 侧视图
  \begin{scope}[shift={(5.6,0)}]
    \draw (0,0) rectangle (4,4);
    \draw[dashed] (1,0) -- (1,4);
    \draw[dashed] (3,0) -- (3,4);
    \draw[<->] (0,4.45) -- (1,4.45) node[midway, above] {$1$};
    \draw[<->] (1,4.45) -- (3,4.45) node[midway, above] {$2$};
    \draw[<->] (3,4.45) -- (4,4.45) node[midway, above] {$1$};
    \draw (0,4) -- (0,4.6);
    \draw (1,4) -- (1,4.6);
    \draw (3,4) -- (3,4.6);
    \draw (4,4) -- (4,4.6);
  \end{scope}
  % 俯视图：圆内接一个边长为 2 的正方形
  \begin{scope}[shift={(2.8,-4.9)}]
    \draw (0,0) circle (2);
    \draw (-1,-1) rectangle (1,1);
  \end{scope}
\end{tikzpicture}
